\chapter{Literature Review}
\label{chap:literature-review}

\section{Review Process}
\label{sec:lr-pra-tools-related}

This chapter reviews contemporary software tools used for probabilistic risk assessment (PRA). The focus is on tools that are currently available and in active use (or at least still distributed and usable today). Legacy or discontinued tools are not reviewed in detail, except where brief historical context is helpful.

The main comparison is summarized in the feature matrix in Table~\ref{tab:pra_tools_overview}. To keep the review consistent and to avoid tool specific descriptions, each software is characterized using the same set of questions. These questions cover the practical, technical and computational aspects of these tools:

\begin{enumerate}
  \item \textbf{Interface:} Does the tool provide a web-based platform that can be accessed through any device? If not, then is it primarily command line based, graphical interface based, or does it support both? Can the CLI be used for batch runs?

  \item \textbf{OS dependency:} Which operating system(s) does the tool depend on? Windows, Linux, macOS, a combination of these 3 or all of them?

  \item \textbf{Licensing:} Is the tool open source, closed source, or free-to-use? Is a license explicitly stated? If a paid license is required, what are the constraints (single user, single machine etc.)?

  \item \textbf{Interoperability:} In what format is the model provided to the tool, and in what format are results produced (e.g., OpenPSA XML, JSInp)? Can the tool import and export models or results for interoperability with other tools?

  \item \textbf{Execution model:} Can the tool execute the tasks in a parallelized way (multi-core), or does it support only serial runs? Does it support distributed, networked execution across multiple machines?

  \item \textbf{Robustness:} If a run fails, becomes stuck, or encounters numerical/modeling errors, does the tool provide error messages and logs that can be used to diagnose the problem?
\end{enumerate}

In addition to reporting tool features, the discussion highlights where capabilities are mature and where gaps remain, especially regarding scalability and interoperability.

\section{Current PRA Tools}
\label{sec:lr-current-pra-tools}

The development of PRA software has progressed through distinct generations, from mainframe based cut set enumerators to contemporary GPU accelerated and decision diagram based solvers.

SAPHIRE \cite{saphire_8, saphire_809} is a long standing PRA suite developed for regulatory work. It is a Windows based desktop application with a GUI that supports graphical fault tree and event tree editing. Historically, it is not a web-based platform. However, recent work describes an extracted solver engine that can run on Linux as well, referred to as SAPHSOLVE \cite{diagnostics_strategies_saphire, refine_saphire_engine, enhancing_saphire_engine}, and explores cloud based solving to improve access to compute resources \cite{advancing_saphire}. SAPHIRE is closed source and its distribution is controlled through agreements and restricted access. SAPHIRE~8 is commonly associated with MAR-D style flat file databases for model and results storage, while SAPHSOLVE uses JSON based inputs/outputs to represent large model logic and solver results. The supported quantifications include MCS generation and probability and frequency quantification, importance measures, uncertainty analysis, and risk metrics such as CDF and LERF. It also supports operational event assessment outputs such as CCDP. The core quantification approach is traditionally cut set based and uses standard approximations, such as min-cut upper bound and rare event approximation alongside Monte Carlo or Latin hypercube sampling for uncertainty. Execution is historically serial for the desktop workflow, while recent plans and prototypes discuss parallel and distributed execution concepts for the extracted solver.

Phoenix Architect \cite{phoenix_architect_v2} is an EPRI managed commercial toolset intended to support a broad PRA workflow, including internal and external event based calculations. It runs on Windows desktop environment with a GUI, while some engines can also be used in batch style workflows through configuration and command line execution options. It is proprietary commercial software with licensing restrictions. Documentation commonly highlights export control constraints and non-disclosure requirements. An educational version exists, but it typically enforces explicit model size limits. The toolset \cite{phoenix_architect_v2_appendices} uses several module specific file formats, for example \texttt{.CAF} for fault trees, \texttt{.ETA} for event trees, and \texttt{.RR} repositories/databases), and commonly exports cut set results (\texttt{.CUT}, \texttt{.RAW}) and report outputs to formats such as Excel or database files. Interoperability is a notable feature: Phoenix Architect is described as supporting exchange with formats such as MAR-D, OpenPSA, and SETS. Quantification capabilities include linked ETs with FTs quantification for CDF and LERF, MCS generation, importance measures (FV/RAW/RRW/Birnbaum), uncertainty analysis, scenario based conditional risk calculations (such as CCDP), and external hazard risk studies (e.g., seismic and fire). Algorithmically, it supports both decision diagram based exact probability computation and cut set based calculations, and uses Monte Carlo or Latin hypercube sampling for uncertainty propagation. In execution mode, it supports both serial and parallel runs, and documentation also mentions remote quantification options where jobs can be executed on separate machines.

CAFTA \cite{cafta_1987, cafta_1989, cafta_10, phoenix_architect_v2_appendices}  is a commercial Windows desktop application for fault tree and event tree modeling and quantification. It is GUI driven and focuses on interactive model building, data management, and reporting. CAFTA is proprietary and typically licensed, with multiple license tiers, for example, professional, educational, and demo versions. CAFTA uses native formats such as \texttt{.CAF} (fault trees) and \texttt{.ETA} (event trees), and commonly stores supporting data in repository/database files (e.g., \texttt{.RR}). Outputs include cut set files (\texttt{.CUT}, \texttt{.RAW}) and report exports (Excel/text). CAFTA supports common PRA quantifications such as MCS generation, top event probability, importance measures, and uncertainty analysis. The cut set generation is typically ZBDD based and cut set probability approximations like MCUB and REA are used for probability calculation, while exact probability calculation is available through BDD based modules. In terms of execution, CAFTA’s core cut set engine is commonly described as serial. Higher performance and parallel execution are often achieved by integrating external solver engines (such as FTREX) rather than relying only on the native engine. For collaboration, CAFTA offers a repository feature that supports multi user teamwork through LAN/server based check in/check out and version history.

FTREX \cite{ftrex_2008, ftrex_2009} is a standalone fault tree solver engine that is commonly used as a backend quantification tool inside larger PRA environments \cite{ftrex_2020}. It is a Windows based application with a CLI interface and is typically commercially licensed with machine bound restrictions. It accepts multiple input formats used in industry workflows, for example, FTAP \texttt{.FTP} and KIRAP \texttt{.KIR}) and can output cut sets in formats such as \texttt{.RAW}, \texttt{.CUT}, and generic text. It can also export logic in formats used for exchange with other ecosystems (e.g., MAR-D files). FTREX supports minimal cut set generation, top event probability computation using cut set approximations (MCUB and REA), importance measures, and multi top quantification workflows that are used for quantifying CDF and LERF. It also supports prevention and path set style calculations for applications such as vital area identification. Algorithmically, FTREX relies on ZBDD structures for efficient MCS generation, and newer versions include BDD based exact probability computation. It also includes delete term approximation options for handling negations in specific settings. FTREX supports parallel execution on a single machine by splitting the problem (for example, by initiators or initiator groups) and solving parts concurrently using a wrapper workflow; this is shared memory parallelism rather than multi node cluster computing.

ZEBRA \cite{zebra_2024} is a recent solver engine focused on accurate quantification of noncoherent fault trees. Public descriptions present it primarily as a computation engine (with no dedicated GUI explicitly described), and the operating system and licensing model are not always clearly stated in publicly available materials. ZEBRA accepts fault tree logic and produces minimal cut sets for coherent models as well as prime implicants (PIs) for noncoherent models. these outputs enable accurate risk metric calculations (e.g., CDF estimation) compared to approaches that rely only on delete term approximations. The core method is based on zero suppressed ternary decision diagrams (ZTDD), which represent positive, negated, and "don’t care" states directly and can therefore encode noncoherent logic. ZEBRA supports both serial and multi thread parallel execution strategies \cite{zebra_2025}. Examples include initiator group partitioning and breadth first parallel solving of gates with reuse of previously computed results across worker threads.

XFTA \cite{xfta_2011, xfta_2020} is a cross platform fault tree calculation engine distributed as freeware (free to use) but not open source. It runs on Windows, Linux, and macOS (64 bit), and it is primarily a CLI tool, with a separate GUI frontend available in some workflows. The software distribution includes explicit restrictions against reverse engineering. XFTA supports OpenPSA XML inputs and also supports a textual, human readable modeling format called S2ML+SBE). Results are typically written as TSV files to support post-processing in spreadsheets and scripts. XFTA supports MCS extraction and ranking, top event probability calculations (including time dependent analysis), a broad set of importance measures, Monte Carlo based sensitivity analysis, and safety integrity level (SIL) metrics such as PFD and PFH. Its methods include a branch-and-test algorithm (SAT inspired) for direct MCS extraction, BDD based exact probability computation, ZBDD representations for compact MCS set handling, and standard cut set probability approximations (REA/MCUB/PUB). XFTA is primarily documented as a serial engine. Performance improvements are mainly achieved through data structure choices, caching, and careful compilation rather than explicit multi threading or distributed execution.

RiskSpectrum \cite{riskspectrum_website, riskspectrum_1990} is a commercial Windows based PRA suite with GUI editors for fault trees and event trees, and modules that support both detailed model development and operational risk monitoring. Its primary modeling environment is desktop based, while the RiskWatcher module uses a client-server architecture for monitoring and viewing risk results in an operational setting (network enabled access to results) \cite{riskspectrum_2010}. RiskSpectrum is proprietary and commercially licensed. Models are typically built using the internal editors, and the tool supports data import and integration workflows, for example, importing plant databases and linking with external fire modeling pipelines. Outputs include minimal cut sets, CDF and LERF results, importance measures and uncertainty results. Scripting support is also described for batch updates and model maintenance. The core quantification is traditionally cut set based, and a BDD solver is available to support exact probability calculations for complex logic. Public sources emphasize the ability to handle very large models, but they often provide limited detail on how much multi core parallelism is used for a single quantification run.

RISKMAN \cite{riskman_2009} is a proprietary commercial software that has a GUI and is commonly presented as part of an industry supported ecosystem, including interfaces to other tools. RISKMAN is described as providing built-in interfaces to other fault tree tools including CAFTA, FTREX, and SAPHIRE to support integrated workflows. Its quantifications emphasize event sequence evaluation, importance measures (including FV/RAW/RRW/Birnbaum and additional fractional measures), and uncertainty analysis using something called "Big Loop" Monte Carlo. It also includes hazard related analysis modules such as fragility based workflows and scenario screening (e.g., seismic and fire screening). Publicly available documentation typically provides limited detail on parallel or distributed execution.

ASTRA \cite{astra_1999, astra_2008, astra_3, astra_3x} is a Windows based GUI tool developed as an institutional fault tree analysis package. It supports importing and exporting the OpenPSA model exchange format and also uses native project and model file formats for storage. Public documentation describes a wide range of quantifications, including unavailability and unreliability analysis (including time dependent evaluation), failure and repair frequency calculations (including PFH measures), expected number of failures over an interval, importance measures, determination of significant minimal cut sets, and uncertainty analysis \cite{astra_2010, astra_2012}. ASTRA is decision diagram based: it uses BDD solving for standard cases, labelled BDD techniques for noncoherent logic, and truncation based methods (e.g., truncated ZBDD and decomposition plus truncation) when full exact representations become too large. Monte Carlo simulation is used for uncertainty propagation. Execution is described mainly as serial.

RiskA is a Windows based PRA tool developed in C\# and \texttt{.NET} ecosystem \cite{riska_2015, riska_5}. It is primarily a desktop GUI based application, and it also has a web-based collaborative modeling feature for fault tree models. Public sources do not clearly state whether RiskA is released under an open source license or a commercial license. For interoperability, RiskA supports a native XML based format and is described as importing OpenPSA MEF XML \cite{riska_2013}, as well as importing model files originating from other industry tools \cite{riska_2016}. Results are commonly exported to Excel and text formats. RiskA supports both qualitative and quantitative fault tree analysis, including MCS generation and prime implicants, top event probability, importance measures (FV/RAW/RRW/Birnbaum), uncertainty and sensitivity analyses, and supporting workflows such as FMEA style reporting. Its core computational method is based on ZBDD \cite{riska_algorithm}, and Monte Carlo simulation is used for uncertainty propagation. The literature does not mention anything about multi threaded or distributed quantification support.

DeRisk \cite{derisk_2018} is a research oriented tool based on Dynamic Uncertain Causality Graphs (DUCG) \cite{ducg_2012}. It is described as a desktop tool with both GUI and CLI interaction modes, and it is implemented in a C\# ecosystem. The licensing model is not clearly specified in the reviewed sources. DeRisk supports classical fault tree quantification including qualitative MCS generation (including disjoint MCSs) and quantitative top event probability, and it supports a broad set of importance measures. A key feature is support for Dynamic Fault Tree behavior through sequential and dependency gates such as PAND, FDEP, SEQ, and SPARE. The working principle described in the literature relies on mapping and probabilistic reasoning within the DUCG representation to capture sequential dependencies without a full state space explosion. Execution is typically described as serial.

\textbf{Hybrid Causal Logic toolset (Trilith, IRIS, and HCLA):}
The Hybrid Causal Logic (HCL) family extends classical FT and ET approaches by integrating event sequence diagrams, fault trees, and Bayesian belief network layers so that causal and socio-technical factors can influence hardware failure logic. In this toolset, IRIS \cite{iris_2007} is commonly described as a desktop GUI modeling environment (primarily Windows), Trilith \cite{trilith_2010} is a commercial Windows desktop application, and HCLA \cite{hcla_web} provides a solver engine that is available both as a cross platform CLI tool and through a web interface (OS agnostic for users). These tools are proprietary/research or commercial offerings rather than open source packages. Inputs are typically XML based model descriptions (often generated from the GUI tools), and outputs include quantitative risk metrics and importance rankings. Quantification capabilities include hybrid solving of ESDs, FTs and BBNs, uncertainty propagation (Monte Carlo, Latin hypercube, and quasi Monte Carlo), and time dependent results such as time-to-failure distributions and time dependent importance measures. The dynamic aspect is usually represented through time models and Bayesian conditioning of probabilities (probabilities change with causal context) rather than through explicit discrete event plant simulation. Public sources generally describe solver execution as serial; the web interface implies server side execution, but details of distributed multi node computation are not typically specified.

OpenPRA \cite{openpra_2023} is an open source web-based platform designed for collaborative PRA modeling and quantification. End users access it through a browser (GUI), while the quantification backend is implemented using containerization techniques, making it suitable for Linux based server environments and portable deployments. It is released under the MIT License. For interoperability, OpenPRA supports OpenPSA MEF XML and also uses JSON based schemas for frontend to backend communication and validation. Outputs can be produced in XML and JSON forms, including cut sets and quantitative results. It supports a broad set of standard PRA quantifications: MCS generation, exact and approximate probability quantification, importance measures, uncertainty analysis via Monte Carlo sampling, event sequence quantification, and CCF analysis. Algorithmically, it relies on engines that implement methods such as BDD and ZBDD, and supports MOCUS based cut set generation as well, and standard cut set approximations (MCUB and REA), alongside Monte Carlo solvers for uncertainty. OpenPRA is designed for distributed execution: it uses a queue based architecture (message broker and worker pool) to distribute jobs across multiple machines and supports concurrent execution of multiple quantifications.

\textbf{SCRAM and its modified versions: }
SCRAM \cite{rakhimov_scram_2022} is a free and open source PRA engine released under GPLv3. It is cross platform (Linux, Windows, macOS) and provides both a CLI interface for headless and batch workflows and a desktop GUI (Qt based) for visualization and interactive use. SCRAM supports OpenPSA MEF XML as its primary input format and produces XML reports. The GUI can export diagrams (e.g., SVG). Supported quantifications include fault tree and event tree quantification, MCS generation and prime implicants (for noncoherent trees), top event probability, importance measures, uncertainty analysis via Monte Carlo sampling, safety integrity level (SIL) metrics, and common cause failure (CCF) analysis using standard CCF models. Methods include MOCUS based cut set generation, BDD based exact probability calculations, ZBDD based set representations, and Monte Carlo sampling for uncertainty. The core solver is generally described as serial for a single quantification run.

SCRAM-cpp \cite{scramcpp_2026} is an enhanced C++ backend engine that builds on the SCRAM ecosystem and is designed for improved performance. It is operated as a CLI solver and is intended to be cross platform, with development and benchmarking often reported on Linux. It continues to use OpenPSA MEF XML as input and output format. A key contribution is explicit shared memory parallelization (e.g., with OpenMP) for computationally heavy tasks such as importance measures and uncertainty sampling, where multi core speedups are reported in benchmark studies. Design discussions also outline directions for distributed memory execution (e.g., MPI) to address problems that exceed single node memory limits, although this is often presented as a roadmap rather than a universally available production feature.

SCRAM++ \cite{scram++_2025} is described as a research platform that wraps multiple quantification algorithms into a unified workflow and selects between them depending on fault tree structure. It is designed to work with OpenPSA MEF inputs and produces standard quantitative outputs (probabilities and cut sets). Its main focus is algorithmic scalability and specialization for applications such as seismic PRA and fragility analysis. A distinguishing method is the use of compressed truth table style formulations to compute exact probabilities efficiently for certain structures such as gates with multiple occurring events, alongside more traditional BDD, ZBDD and MOCUS methods. Public descriptions emphasize serial performance improvements and hybrid method selection more than explicit parallel or distributed execution.

mcSCRAM \cite{arjun_mcscram} is a high performance Monte Carlo engine developed in the SCRAM ecosystem and designed to run efficiently on multi core CPUs and GPUs using SYCL based backend. It is a CLI tool and does not provide a dedicated web interface on its own. mcSCRAM is released under the AGPLv3 license, which has important implications for network deployment (source availability requirements when used as a service). It supports OpenPSA MEF inputs and produces probability estimates and statistical outputs consistent with the SCRAM style reporting approach. Unlike deterministic cut set or BDD solvers, mcSCRAM focuses on data parallel Monte Carlo: it compiles the model into a probabilistic circuit (a DAG representation of the logic), evaluates many Monte Carlo trials at once using bit-packed Boolean vectors, and uses counter based Philox random number generation to support massive parallel sampling without synchronization overhead. This design provides strong parallel performance on a single heterogeneous node (GPU and CPU). Public descriptions generally do not claim multi node distributed execution across a cluster, and the tool is described as being in an early/alpha stage.

\section{Existing Distributed Computing Approaches for PRA}
\label{sec:lr-parallel-distributed}

Across the reviewed tools, most PRA quantification still assumes a single machine execution model. In many commercial desktop environments, the default workflow is serial execution on a Windows workstation, with scalability achieved mainly through solver efficiency, approximation settings, or splitting the study into multiple independent runs. Parallelism is more common in (i) uncertainty analyses that require many repeated samples (e.g., SCRAM-cpp), (ii) batch evaluation of many sequences or many top events (e.g., Phoenix Architect, FTREX), and (iii) Monte Carlo simulation engines that can spread trials across CPU cores or GPU threads (e.g., mcSCRAM).

A smaller set of tools describe stronger forms of parallel execution. Some solver engines parallelize by partitioning a large model into smaller independent subproblems and solving them concurrently on a multi core machine (e.g., ZEBRA). Some open source engines explicitly parallelize internal loops for importance and uncertainty calculations using shared memory programming models (e.g., SCRAM-cpp). Monte Carlo based tools can achieve substantial speedups on GPUs (e.g., mcSCRAM), but they often focus on saturating a single heterogeneous node rather than scaling across multiple networked machines.

Distributed, networked execution where multiple worker machines are coordinated through a queue and a central scheduler is described much less often in public documentation. Among the reviewed tools, OpenPRA provides a clear end-to-end distributed design: it uses a message broker, a scalable worker pool, task decomposition, and operational fault handling mechanisms. A few commercial toolchains mention the ability to run quantification on remote machines (e.g., Phoenix Architect), but public sources typically provide fewer details about scheduling semantics and multi tenant queueing. Overall, the literature shows growing interest in web access and scalable computation, but openly specified distributed queueing infrastructures for heterogeneous PRA workloads are still uncommon.

\newcolumntype{L}[1]{>{\raggedright\arraybackslash\hsize=#1\hsize}X}
\newcolumntype{C}[1]{>{\centering\arraybackslash\hsize=#1\hsize}X}
\begin{landscape}
\begin{table}
\centering
\begin{threeparttable}
\scriptsize
\setlength{\tabcolsep}{4pt}
\caption{Feature overview of reviewed \acrshort{pra} tools.}
\label{tab:pra_tools_overview}

\begin{tabularx}{\textwidth}{@{} L{1.20} L{1.30} L{0.60} L{1.55} L{0.85} C{0.50} @{}}
\toprule
\textbf{Tool} &
\textbf{Platform/OS} &
\textbf{User} &
\textbf{Model Exchange Format} &
\textbf{License} &
\textbf{Processing Mode} \\
\midrule

\acrshort{saphire} &
Desktop/Windows &
Single &
MAR-D + JSInp &
Proprietary &
S \\

Phoenix Architect &
Desktop/Windows &
Multi\tnote{*} &
.CAF, .ETA, .RR, .CUT, .RAW + MAR-D + OpenPSA + SETS &
Commercial proprietary &
S/P/D\tnote{$\dagger$} \\

CAFTA &
Desktop/Windows &
Multi\tnote{*} &
.CAF, .ETA, .RR, .CUT, .RAW + SETS + FTAP &
Commercial proprietary &
S \\

\acrshort{ftrex} &
Desktop/Windows &
Single &
FTAP/.FTP + KIRAP/.KIR + .RAW, .CUT, .TXT + MAR-D &
Commercial proprietary &
P \\

ZEBRA &
Solver engine/OS n/s &
Single &
Unspecified &
Unspecified &
S/P \\

XFTA &
Desktop/Cross platform &
Single &
OpenPSA + S2ML+SBE + TSV &
Freeware proprietary &
S \\

RiskSpectrum &
Desktop/Windows + client-server &
Single\tnote{*} &
Unspecified &
Commercial proprietary &
S \\

RISKMAN &
Desktop/Windows &
Single &
CAFTA + FTREX + SAPHIRE &
Commercial proprietary &
S \\

ASTRA &
Desktop/Windows &
Single &
OpenPSA + .tfx, .adf &
Proprietary &
S \\

RiskA &
Desktop/Windows + client-server &
Multi &
XML + OpenPSA + CAFTA and RiskSpectrum formats + Rich text + Excel, text, Word &
Proprietary &
S \\

DeRisk &
Desktop/OS n/s &
Single &
Unspecified &
Proprietary &
S \\

Trilith &
Desktop/Windows &
Single &
XML export &
Commercial proprietary &
S \\

\acrshort{hcla} &
CLI/Cross platform + client-server &
Multi &
JSON &
Proprietary &
S \\

OpenPRA &
Client-server &
Multi &
OpenPSA + OpenPRA JSON &
MIT + GPLv3 &
D \\

SCRAM &
Desktop/Cross platform + Solver engine &
Single &
OpenPSA + SVG &
GPLv3 &
S \\

SCRAM-cpp &
Solver engine/Cross platform &
Single &
OpenPSA &
GPLv3 &
P \\

SCRAM++ &
Solver engine/Cross platform &
Single &
OpenPSA &
Proprietary &
S \\

mcSCRAM &
Solver engine/Cross platform &
Single &
OpenPSA &
AGPLv3 &
P \\

\bottomrule
\end{tabularx}

\begin{tablenotes}
  \item[*] ``Multi-user'' here refers to built-in support for team workflows (e.g., repositories or web access), not necessarily simultaneous editing of the same model without coordination.
  \item[$\dagger$] Phoenix Architect documentation mentions remote quantification options. Public sources do not describe full scheduling and fault tolerance semantics.
  \item[] Processing types: S = serial; P = parallel on a single machine (shared memory / GPU); D = distributed execution across multiple machines.
  \item[] n/s = not specified clearly in the reviewed public documentation.
\end{tablenotes}

\end{threeparttable}
\end{table}
\end{landscape}

\section{Limitations of Current PRA Tools}
\label{sec:lr-limitations}

The current software ecosystem of PRA exhibits recurring limitations that affect scalability, interoperability, and operational reliability. While some tools may partially address these problems, they are either proprietary or not yet fully developed. The following constraints are broadly observed across widely used desktop applications, command line quantifiers, and web-based platforms:

\begin{itemize}
  % \item External hazard models (e.g., fire, seismic) frequently contain tens of thousands of basic events, causing memory explosion and extended runtimes (hours or days) during minimal cut set (MCS) manipulation and decision diagram construction (BDD, ZBDD), particularly with noncoherent logic and deeply nested expansions.

  \item Operating system dependencies (often Windows centered) and local GUI runtimes hinder cross platform reproducibility and complicate execution in controlled compute environments (e.g., Linux based HPC clusters).

  \item Proprietary commercial tools impede external optimization for multi core, GPU, or cluster environments, limit independent verification, and restrict community driven improvements.
  
  % while open source alternatives often lack robust multi tenant service infrastructure and operational tooling expected in production settings.

  \item Most tools support only serial or single node parallel execution. Distributed submission interfaces are comparatively rare and, when present, are commonly tied to private script based schedulers rather than modern network accessible REST APIs with explicit resource governance.

  \item Current tools generally lack support for high volumes of simultaneous job requests with asynchronous execution, resulting in unpredictable turnaround times and inadequate priority handling for heterogeneous workloads (e.g., point estimate runs take less time compared to uncertainty quantification, importance measures and sensitivity studies).

  \item The absence of a unified platform that can host multiple solvers forces users to buy licenses for multiple solvers perform error prone conversions across different proprietary formats, impeding solver interoperability and making it difficult to compare results across quantification engines under controlled conditions.

  \item Sensitive nuclear plant data frequently necessitates on premise deployment, yet many tools lack built in identity and access management (IAM) integration and role based access controls (RBAC).

  % \item Per host licensing policies constrain scaling across on-premise cluster nodes, often limiting deployments to single node configurations, thereby reducing achievable throughput for large model batches.

  \item Most quantifiers lack native GPU support and heterogeneous resource orchestration, and any available thread level parallelism is rarely coordinated by a multi tenant scheduler.
  
  % This can lead to rapid saturation of available resources under concurrent usage.

  % \item Many quantifiers tightly couple parsing, compilation, quantification, and post processing in a single monolithic run, hindering scalable pre processing (e.g., parallel transformations of event tree sequences into individual fault trees), incremental compilation, and reuse of intermediate artifacts.

  % \item Standardized reporting of performance and service quality metrics (e.g., latency, throughput, queue wait time, solver utilization, error rates) is uncommon. In the absence of such metrics, it is difficult to set technical requirements for web-based quantification services under multi user demand or to diagnose capacity bottlenecks.

  \item When large runs fail, stall, or encounter memory errors, tools typically print error messages but often do not indicate which specific part of the model (e.g., a particular event sequence) caused the problem or where the run became stuck. Logs are usually limited to overall run status (success/failure) or generic error messages, which makes triage and root cause analysis difficult. In contrast, a distributed networked execution system can isolate failures to specific sub jobs (e.g., sequence level tasks), retain per task logs, and provide actionable diagnostics.
\end{itemize}

Collectively, these limitations highlight the critical need for automated and portable deployment options, standardized interfaces, and scalable execution strategies that maintain data governance while accelerating processing times for large, realistic PRA models. They also motivate PRA execution platforms that treat quantification as an observable service: one that can schedule work, isolate faults, retain structured logs, and provide reproducible execution records across heterogeneous solvers and workloads.

\section{Summary}
\label{sec:lr-summary}

The PRA software ecosystem spans legacy cut set enumerators, mature desktop suites, command line quantifiers, causal modeling platforms, and emerging web enabled services. Each category offers distinctive strengths in areas such as comprehensive modeling environments, advanced quantification algorithms, or improved accessibility. However, the literature and current practice reveal some persistent limitations. These include OS specific deployments, constrained parallelism, limited support for distributed and asynchronous execution, proprietary formats that inhibit solver interoperability, weak service level observability, and incomplete security features for on-premise deployments. In addition, robustness remains a practical operational gap: when large batch computations fail or hang, available tooling often provides only coarse error output that does not pinpoint the responsible model component (e.g., an individual event sequence), increasing the time and effort required for debugging and rerunning large studies.

Prior works demonstrate important algorithmic improvements (e.g., decision diagram techniques) and introduces web-based mechanisms for remote execution. Nonetheless, the literature lacks comprehensive descriptions of a fully specified, solver agnostic, publicly available queueing and execution system designed to support heterogeneous PRA workloads under multi user demand, with explicit capabilities for prioritization, resource governance, structured logging, and failure isolation. The scalability, interoperability, robustness, and governance challenges documented in this review motivate the development of a systematic, high performance distributed computational solution, which forms the focus of subsequent chapters.
